\section{Proofs}

\subsubsection*{Proof of Proposition \ref{prop_pay_to_post}}

\begin{proof}
A post idea $(q, \epsilon)$ is posted if and only if $\epsilon > \kappa - \ln V(q)$. Define $\pi(q,\kappa)$ as the probability (over $\epsilon$) that a post idea of quality $q$ is posted when the log cost is $\kappa$:
\begin{align*}
\pi(q,\kappa) &= P(\epsilon > c - \ln V(q)) 
\end{align*}
$\pi(q,\kappa)$ is the survival function of $\epsilon|q$, evaluated at $\kappa - \ln V(q)$: $\pi(q,\kappa) = S_{\epsilon|q}(\kappa - \ln V(q) | q)$.
\begin{enumerate}[label=(\roman*)]
\item The unconditional probability of posting is:
\begin{align*}
P(\text{post}) = \int S_{\epsilon|q}(\kappa- \ln V(q)|q)f_q(q)dq
\end{align*}
Since $S_{\epsilon|q}(\cdot|q)$ is a survival function, its first derivative is negative everywhere. An increase in $c$ will therefore reduce the value of $S_{\epsilon|q}(\kappa - \ln V(q))$ at every $q$. Since that is what we are integrating over to get $P(\text{post})$, this strictly reduces $P(\text{post})$. The expected number of posts per unit time is $\gamma P(\text{post})$, so an increase in $c$ strictly reduces the expected number of posts.

\item Take any $c^\prime > c$ with corresponding log costs $\kappa^\prime > \kappa$. The set of ideas that were posted at log cost $\kappa$ can be partitioned into two groups:
\begin{itemize}
\item \textbf{Surviving posts.} Ideas that are still posted at log cost $\kappa^\prime$, i.e. $\ln V(q) + \epsilon > \kappa^\prime$
\item \textbf{Marginal posts.} Ideas that are posted at $\kappa$ but not at $\kappa^\prime$, i.e. $\kappa < \ln V(q) + \epsilon < \kappa^\prime$.
\end{itemize}
By the Law of Iterated Expectations,
\begin{align*}
E[q | \text{post at } \kappa] = \alpha E[q | \text{surviving}] + (1-\alpha) E[q | \text{marginal}]
\end{align*}
where $\alpha = P(\text{surviving} | \text{post at } \kappa)$. Noting that $E[q|\text{post at } \kappa^\prime] = E[q|\text{surviving}]$, we rearrange the above equation to obtain:
\begin{align*}
E[q|\text{post at } \kappa^\prime] - E[q|\text{post at } \kappa] = (1-\alpha) \left( E[q|\text{surviving}] - E[q|\text{marginal}\right)
\end{align*}
Thus, to show that expected post quality is increasing in $c$, it suffices to show that the expected quality of the surviving posts is always higher than the expected quality of the marginal posts. 

To demonstrate this, it is sufficient to establish that for any $q_2 > q_1$, the ratio $\pi(q_2,\kappa)/\pi(q_1,\kappa)$ is strictly increasing in $\kappa$. Intuitively, this says that as costs rise, higher quality posts retain a larger share of the posting probability relative to lower quality ideas. This would imply that the marginal posts are tilted toward lower quality than the surviving posts.

We write:
\begin{align*}
\frac{\pi(q_2,\kappa)}{\pi(q_1,\kappa)} = \frac{S_{\epsilon|q}(\kappa-\ln V(q_2)|q_2)}{S_{\epsilon|q}(\kappa - \ln V(q_1)|q_1)} \label{eq_prob_ratio}
\end{align*}
Define the hazard rate as $h(\epsilon|q) = f_{\epsilon|q}(\epsilon|q) / S_{\epsilon|q}(\epsilon|q)$. Differentiating $\pi(q_2,\kappa)/\pi(q_1,\kappa)$ with respect to $\kappa$ shows us that the ratio is increasing in $\kappa$ if and only if:
\begin{align*}
h(\kappa - \ln V(q_1) | q_1) > h(\kappa - \ln V(q_2) | q_2)
\end{align*}
The proof now proceeds in two parts. 

\paragraph{Claim 1.} First, we show that $h(\epsilon | q_1) \geq h(\epsilon | q_2)$ for any $q_2 > q_1$. By the MLRP, for any $t > \epsilon$, we have that $\ell(t; q_1, q_2) \geq \ell(\epsilon; q_1, q_2)$, which rearranges to:
\begin{align*}
f_{\epsilon|q}(t|q_2) f_{\epsilon|q}(\epsilon|q_1) \geq f_{\epsilon|q}(\epsilon|q_2) f_{\epsilon|q}(t|q_1)
\end{align*}
Integrating both sides over $t$ from $\epsilon$ to $\infty$ gives us:
\begin{align*}
S_{\epsilon|q}(\epsilon|q_2) f_{\epsilon|q}(\epsilon|q_1) \geq S_{\epsilon|q}(\epsilon|q_1) f_{\epsilon|q}(\epsilon|q_2)
\end{align*}
which rearranges to $h(\epsilon|q_1) \geq h(\epsilon | q_2)$ as was to be demonstrated.

\paragraph{Claim 2.} Now, we show that $h(\epsilon|q)$ is increasing in $\epsilon$ for each $q$. Take the derivative of $h(\epsilon|q)$ with respect to $\epsilon$. The sign is equal to the sign of:
\begin{align*}
f_{\epsilon|q}(\epsilon|q)^2 + f^\prime_{\epsilon|q}(\epsilon|q) S_{\epsilon|q}(\epsilon|q)
\end{align*}
When $f^\prime_{\epsilon|q}(\epsilon|q) \geq 0$, this is immediately positive. When $f_{\epsilon|q}^\prime(\epsilon|q)<0$, we rely on the log-concavity property. Since $f_{\epsilon|q}(\cdot)$ is log-concave, the following holds for every $t > \epsilon$:
\begin{align*}
\ln f_{\epsilon|q} (t|q) < \ln f_{\epsilon|q}(\epsilon|q) + (t-\epsilon) \frac{f_{\epsilon|q}^\prime(\epsilon|q)}{f_{\epsilon|q}(\epsilon|q)}
\end{align*}
Exponentiating both sides:
\begin{align*}
f_{\epsilon|q} (t|q) < f_{\epsilon|q}(\epsilon|q) \exp\left[ (t-\epsilon) \frac{f_{\epsilon|q}^\prime(\epsilon|q)}{f_{\epsilon|q}(\epsilon|q)} \right]
\end{align*}
Now integrating over $t$ from $\epsilon$ to $\infty$:
\begin{align*}
S_{\epsilon|q} (\epsilon | q) &< f_{\epsilon|q}(\epsilon|q) \int_{\epsilon}^\infty \exp\left[ (t-\epsilon) \frac{f_{\epsilon|q}^\prime(\epsilon|q)}{f_{\epsilon|q}(\epsilon|q)} \right] dt \\
& = - \frac{f_{\epsilon|q}(\epsilon|q)^2}{f^\prime_{\epsilon|q}(\epsilon|q)}
\end{align*}
Since $f^\prime_{\epsilon|q}(\epsilon|q) < 0$, this can be rearranged to:
\begin{align*}
f_{\epsilon|q}(\epsilon|q)^2 + f^\prime_{\epsilon|q}(\epsilon|q) S_{\epsilon|q}(\epsilon|q) > 0
\end{align*}
which was to be demonstrated. Thus, $h(\epsilon|q)$ is strictly increasing in $\epsilon$ for all $q$.

By Assumption \ref{asm_monotonicity}, $\kappa - \ln V(q_1) > \kappa - \ln V(q_2)$. By Claim 2, this implies that $h(\kappa - \ln V(q_1)|q_1) > h(\kappa-\ln V(q_2)|q_1)$. By Claim 1, $h(\kappa-\ln V(q_2) | q_1) \geq h(\kappa-\ln V(q_2) | q_2)$. Together, we have $h(\kappa - \ln V(q_1)|q_1) > h(\kappa - \ln V(q_2)|q_2)$, which was to be demonstrated. 
\end{enumerate}
\end{proof}


